\begin{aufgabe}[Gedächtnisprotokoll, Aufgabe 4]
\begin{itemize}
  \item $[10, 8, 3]_{11}$ Code und die folgenden Wörter wurden empfangen
  \begin{itemize}
    \item $r=115000711$
    \item $r=115000733$
  \end{itemize}
\end{itemize}
% UNSICHER: Empfangene Wörter haben in der Quelle nur 9 Stellen, obwohl der Code Länge 10 hat; im Text als r_1 und r_2 bezeichnet.
\begin{enumerate}
  \item[a)] Stellen Sie die Kontrollmatrix auf
  \item[b)] Berechnen Sie Syndrom von $r_1$ und $r_2$
  \item[c)] Welches der beiden Wörter ist nicht dekodierbar und wieso?
  \item[d)] Korrigieren Sie das entsprechende Wort, das dekodierbar ist (also bestimmen Sie $c$)
  \item[c)] Zeigen Sie, dass das korrigierte Wort tatsächlich ein Codewort ist
  % UNSICHER: Teilaufgabe in der Quelle erneut als "c)" bezeichnet (vermutlich e)).
\end{enumerate}
% Punktzahl: nicht angegeben.
\end{aufgabe}
